\section{第四章：数学物理方程}

\begin{frame}{第四章：从物理到偏微分方程}
  \begin{itemize}
    \item 选择待求量：位移、温度或势函数。
    \item 取一个微小区域，列写受力、通量或守恒关系。
    \item 令区域尺度趋于零，得到局部微分方程。
    \item 配合初始条件和边界条件，确定物理解。
  \end{itemize}
  \begin{block}{建模原则}
    方程不仅编码规律，也编码假设；改变介质、尺度或近似，模型可能必须改变。
  \end{block}
\end{frame}

\begin{frame}{波动方程：张力与惯性的平衡}
  设 $u(x,t)$ 是绷紧细弦的横向位移，线密度为 $\rho$、张力为 $T$。小振幅近似下，微段的横向合力为
  \[
    T\bigl[u_x(x+\Delta x,t)-u_x(x,t)\bigr].
  \]
  由牛顿第二定律并令 $\Delta x\to0$，得到
  \[
    \rho u_{tt}=Tu_{xx},
    \qquad
    u_{tt}-a^2u_{xx}=0,
    \qquad a^2=\frac{T}{\rho}.
  \]
  \begin{itemize}
    \item $u_{tt}$ 表示惯性，$u_{xx}$ 表示弦的弯曲程度。
    \item 双曲型方程描述有限速度传播，必须给定初始位移和初始速度。
  \end{itemize}
\end{frame}

\begin{frame}{热方程：通量与能量守恒}
  Fourier 定律指出热流方向与温度梯度相反：
  \[
    \mathbf{q}=-k\nabla u.
  \]
  对任意体积使用能量守恒和散度定理，得到
  \[
    \rho c\,u_t=\nabla\cdot(k\nabla u).
  \]
  对均匀各向同性介质，
  \[
    u_t=\kappa\nabla^2u,
    \qquad \kappa=\frac{k}{\rho c}.
  \]
  \begin{block}{物理图像}
    温度高峰会向周围扩散并逐渐平滑；这与波动方程的“传播波形”有本质区别。
  \end{block}
\end{frame}

\begin{frame}{Laplace 与 Poisson 方程}
  当系统达到静态平衡，时间导数消失，常见的椭圆型方程为
  \[
    \nabla^2u=0\quad\text{(Laplace 方程)},
    \qquad
    \nabla^2u=f\quad\text{(Poisson 方程)}.
  \]
  \begin{columns}[T]
    \column{0.48\textwidth}
    \begin{block}{无源区域}
      电势、稳态温度等满足 Laplace 方程。内部值由边界值强烈约束。
    \end{block}
    \column{0.48\textwidth}
    \begin{block}{有源区域}
      $f$ 表示电荷、热源或外力等分布式来源。
    \end{block}
  \end{columns}
  \[
    \nabla^2u=u_{xx}+u_{yy}+u_{zz}.
  \]
\end{frame}

\begin{frame}{分类决定应问什么问题}
  对二阶线性 PDE，典型分类与物理行为如下：
  \begin{center}
    \begin{tabular}{lll}
      类型 & 代表方程 & 主要物理特征\\ \hline
      双曲型 & $u_{tt}-a^2u_{xx}=0$ & 波以有限速度传播\\
      抛物型 & $u_t-\kappa u_{xx}=0$ & 扩散、平滑与耗散\\
      椭圆型 & $\nabla^2u=0$ & 静态平衡与边界控制
    \end{tabular}
  \end{center}
  \begin{alertblock}{解题前先问}
    未知量是什么？区域在哪里？给定的是初始条件、边界条件，还是两者都有？这些信息决定问题是否适定。
  \end{alertblock}
\end{frame}

\begin{frame}{分离变量：把 PDE 化为 ODE}
  对边界条件 $u(0,t)=u(\ell,t)=0$ 的热方程，尝试
  \[
    u(x,t)=X(x)T(t).
  \]
  代入 $u_t=\kappa u_{xx}$ 并分离变量：
  \[
    \frac{T'}{\kappa T}=\frac{X''}{X}=-\left(\frac{n\pi}{\ell}\right)^2.
  \]
  因此
  \[
    u_n(x,t)=A_n\sin\left(\frac{n\pi x}{\ell}\right)
    e^{-\kappa n^2\pi^2t/\ell^2}.
  \]
  \begin{block}{关键步骤}
    边界条件选择允许的空间模态；线性性允许把所有模态叠加，初始条件决定各模态系数。
  \end{block}
\end{frame}
